% sec:dg-manifolds — Manifolds, Tangent Spaces, and One-Forms
\section{Manifolds, Tangent Spaces, and One-Forms}
\label{sec:dg-manifolds}

Try covering a sphere with a single coordinate grid.  Near any point, latitude
and longitude work fine --- but at the poles, longitude degenerates to a point and
the coordinate lines pile up.  No single set of coordinates can describe the
entire sphere smoothly.  The same problem arises for any curved spacetime: the
manifold might have topology that no single coordinate chart can cover.  The
resolution is to work with patches --- overlapping coordinate systems, each
covering part of the manifold, with smooth transition rules where they overlap.

\paragraph{Smooth manifolds.}

A \emph{smooth manifold} $\mathcal{M}$ of dimension $n$ is a topological space
that is Hausdorff, second-countable, and locally homeomorphic to
$\mathbb{R}^n$, equipped with a maximal collection of compatible coordinate
charts.  A \emph{chart} $(U, \varphi)$ consists of an open set $U \subset
\mathcal{M}$ and a homeomorphism $\varphi : U \to \varphi(U) \subset
\mathbb{R}^n$ that assigns coordinates to each point.  Where two charts
$(U, \varphi)$ and $(V, \psi)$ overlap, a natural question arises: how do the
coordinates from one chart relate to the coordinates from the other?  The answer
is the \emph{transition map}
%
\begin{equation}
  \psi \circ \varphi^{-1} : \varphi(U \cap V) \to \psi(U \cap V) \,,
  \label{eq:transition-map}
\end{equation}
%
which is a map from one open subset of $\mathbb{R}^n$ to another --- ordinary
multivariable calculus applies here, since both sides live in $\mathbb{R}^n$.
The manifold is \emph{smooth} if every such transition map is infinitely
differentiable ($C^\infty$).  An \emph{atlas} is a collection of charts whose
domains cover $\mathcal{M}$; the \emph{maximal} atlas includes every chart
compatible with the existing ones, ensuring that the smooth structure doesn't
depend on a particular choice of initial charts.
\histnote{The modern definition of a manifold was formalised by Whitney in the
1930s, though the idea of working with local coordinate patches goes back to
Riemann's 1854 Habilitationsschrift.}

\begin{figure}[htbp]
  \centering
  % manifold-charts.tex — Two overlapping coordinate charts on S²
% Inputable TikZ fragment for fig:manifold-charts
% Requires: tikz with calc, arrows.meta, positioning (loaded by ehbook.cls)
%
\begin{tikzpicture}[
    >={Stealth[length=5pt]},
    font=\footnotesize,
    map arrow/.style={->, thick, EHMarginGray},
  ]

  % ── Left R² patch (gold) ──────────────────────────────────────────
  \fill[EHAccretionGold!15, draw=EHMarginGray, thick, rounded corners=2pt]
    (-0.9,-1.15) rectangle (0.9,1.15);
  \node[font=\small] at (0,0) {$\varphi(U \cap V)$};

  % ── Centre sphere S² ──────────────────────────────────────────────
  \begin{scope}[shift={(3.5,0)}]
    % Southern cap (blue)
    \begin{scope}
      \clip (0,0) circle (1.1);
      \fill[EHJetBlue!20] (-1.2,-1.2) rectangle (1.2,-0.3);
    \end{scope}
    % Overlap band
    \begin{scope}
      \clip (0,0) circle (1.1);
      \fill[EHAccretionGold!50!EHJetBlue, fill opacity=0.15]
        (-1.2,-0.3) rectangle (1.2,0.3);
    \end{scope}
    % Northern cap (gold)
    \begin{scope}
      \clip (0,0) circle (1.1);
      \fill[EHAccretionGold!20] (-1.2,0.3) rectangle (1.2,1.2);
    \end{scope}

    % Outline
    \draw[EHMarginGray, thick] (0,0) circle (1.1);

    % Overlap boundary lines
    \pgfmathsetmacro{\xov}{sqrt(1.1*1.1 - 0.3*0.3)}
    \draw[EHMarginGray, thin, densely dashed] (-\xov, 0.3) -- (\xov, 0.3);
    \draw[EHMarginGray, thin, densely dashed] (-\xov,-0.3) -- (\xov,-0.3);

    % Labels
    \node[font=\small, text=EHAccretionGold!80!black] at (0, 0.72) {$U$};
    \node[font=\small, text=EHJetBlue!80!black]       at (0,-0.72) {$V$};
    \node[fill=white, inner sep=1pt, font=\scriptsize] at (0, 0)   {$U \!\cap\! V$};
    \node[font=\small, above=2pt]                      at (0, 1.1) {$S^2$};
  \end{scope}

  % ── Right R² patch (blue) ─────────────────────────────────────────
  \fill[EHJetBlue!15, draw=EHMarginGray, thick, rounded corners=2pt]
    (6.1,-1.15) rectangle (7.9,1.15);
  \node[font=\small] at (7.0,0) {$\psi(U \cap V)$};

  % ── Map arrows ────────────────────────────────────────────────────

  % φ: sphere overlap → left R² patch
  \draw[map arrow]
    (3.5-1.1, 0.1) to[out=155, in=25]
    node[above, pos=0.5] {$\varphi$}
    (0.9, 0.3);

  % ψ: sphere overlap → right R² patch
  \draw[map arrow]
    (3.5+1.1, 0.1) to[out=25, in=155]
    node[above, pos=0.5] {$\psi$}
    (6.1, 0.3);

  % Transition map: left patch → right patch
  \draw[map arrow]
    (0,-1.55) -- node[below=1pt] {$\psi \circ \varphi^{-1}$} (7.0,-1.55);

\end{tikzpicture}

  \caption{Two overlapping charts on the two-sphere $S^2$.  The chart
    $(U, \varphi)$ (gold) covers the northern hemisphere; $(V, \psi)$ (blue)
    covers the southern hemisphere.  In the overlap region (shaded band near
    the equator), the transition map $\psi \circ \varphi^{-1}$ translates
    between the two coordinate representations.  The manifold is smooth
    because this transition map is $C^\infty$.  No single chart covers the
    entire sphere --- the coordinate singularity at the poles forces us to
    use at least two patches.}
  \label{fig:manifold-charts}
\end{figure}

For spacetime, $n = 4$: each chart assigns four coordinates $x^\mu =
(x^0, x^1, x^2, x^3)$ to events in a patch, with the Einstein summation
convention applying to repeated upper--lower index pairs as in
\cref{ch:special-relativity}.  The coordinates used there --- Cartesian
$(t, x, y, z)$ --- constitute a single chart that happens to cover all of
Minkowski spacetime.  For Schwarzschild or Kerr spacetimes, a single chart is
insufficient: the Boyer--Lindquist coordinates break down at the event horizon,
and different patches (such as Eddington--Finkelstein coordinates) are needed to
cover the full manifold.
\xrefnote{Coordinate singularities at the horizon are discussed in
\cref{ch:schwarzschild}.}

\paragraph{Tangent vectors as directional derivatives.}

In Euclidean space, a vector is an arrow: a displacement with magnitude and
direction.  On a curved manifold, this picture fails --- there is no ambient
space to embed arrows in, and ``displacing along a direction'' doesn't produce
a unique result on a curved surface.  The modern resolution requires a genuine
shift in perspective: we define tangent vectors not as arrows but as
\emph{directional derivative operators}.  This is one of the conceptual leaps
in differential geometry, and it deserves a careful look.

At a point $p \in \mathcal{M}$, a \emph{tangent vector} $V$ is a linear map
from smooth functions (more precisely, from germs of smooth functions at $p$)
to real numbers that satisfies the Leibniz rule:
%
\begin{equation}
  V(fg) = f(p)\,V(g) + g(p)\,V(f) \,.
  \label{eq:tangent-vector-leibniz}
\end{equation}
%
The germ condition means that $V$ depends only on the behaviour of $f$ in an
arbitrarily small neighbourhood of $p$ --- if two functions agree near $p$,
$V$ gives the same result on both.  Concretely, $V$ is the familiar
directional derivative, promoted to a definition: a tangent vector tells you
``if I move in this direction, how fast does the function $f$ change?''  The
collection of all tangent vectors at $p$ forms a vector space called the
\emph{tangent space} $T_p\mathcal{M}$.  For a four-dimensional manifold, the
tangent space is four-dimensional.
\physnote{A particle's four-velocity $\fourvel$ is a tangent vector to its
worldline --- the derivative of the worldline's coordinate functions with
respect to proper time $\tau$, normalised so that $g_{\mu\nu}\,u^\mu\,u^\nu
= -1$ for massive particles.}

Given coordinates $x^\mu$ in a chart containing $p$, the coordinate basis
vectors for $T_p\mathcal{M}$ are the partial derivative operators:
%
\begin{equation}
  \mathbf{e}_{(\mu)} = \frac{\partial}{\partial x^\mu}
    \equiv \pd{\mu} \,.
  \label{eq:coordinate-basis}
\end{equation}
%
Any tangent vector can be expanded in this basis as $V = V^\mu\,\pd{\mu}$,
where the $V^\mu$ are the contravariant components we used throughout
\cref{ch:special-relativity}.  The basis depends on the coordinate chart --- a
different chart gives different basis vectors --- but the tangent vector $V$
itself is a geometric object, independent of coordinates.

\paragraph{One-forms and the cotangent space.}

For every tangent space $T_p\mathcal{M}$, there is a dual space: the
\emph{cotangent space} $T^*_p\mathcal{M}$, consisting of linear maps from
tangent vectors to real numbers.  An element $\omega \in T^*_p\mathcal{M}$ is
called a \emph{one-form} or \emph{covector}.

The coordinate basis for the cotangent space is the set of differentials
$\{\d x^\mu\}$, defined by their action on the coordinate basis vectors:
%
\begin{equation}
  \d x^\mu(\pd{\nu}) = \delta^\mu{}_\nu \,.
  \label{eq:dual-basis}
\end{equation}
%
This is the defining property of a dual basis: each basis one-form picks out
the corresponding coordinate component and ignores the rest.  An arbitrary
one-form can be expanded as $\omega = \omega_\mu\,\d x^\mu$, where the
$\omega_\mu$ are the covariant components from \cref{sec:sr-minkowski}.

The pairing of a one-form $\omega$ with a tangent vector $V$ is the natural
contraction:
%
\begin{equation}
  \omega(V) = \omega_\mu\, V^\mu \,.
  \label{eq:covector-pairing}
\end{equation}
%
This is the same index contraction we've been using since
\cref{ch:special-relativity}, but now with a precise geometric meaning: the
one-form $\omega$ eats a tangent vector $V$ and produces a number.  No metric
is required for this operation --- it is purely algebraic.  The metric enters
only when we want to convert between vectors and one-forms (raising and
lowering indices), which is the subject of \cref{sec:dg-metric}.
\warnnote{The pairing $\omega(V) = \omega_\mu V^\mu$ does not involve the
metric.  This is a common source of confusion: contraction between upper and
lower indices is an algebraic operation, not a geometric one.  The metric is
needed only for \emph{raising} or \emph{lowering} indices.}

\paragraph{Coordinate transformations.}

Under a change of coordinates $x^\mu \to x'^\mu$, the basis vectors and basis
one-forms transform inversely.  For the tangent vectors:
%
\begin{equation}
  \pd{\mu'} = \frac{\partial x^\nu}{\partial x'^\mu}\,\pd{\nu} \,.
  \label{eq:basis-transform}
\end{equation}
%
This says that the new basis vector $\pd{\mu'}$ is a linear combination of the
old basis vectors, with the coefficients given by the inverse Jacobian matrix
$\partial x^\nu / \partial x'^\mu$.  For the basis one-forms, the
transformation goes the other way:
%
\begin{equation}
  \d x'^\mu = \frac{\partial x'^\mu}{\partial x^\nu}\,\d x^\nu \,.
  \label{eq:cobasis-transform}
\end{equation}
%
The coefficients here are the Jacobian matrix $\partial x'^\mu / \partial
x^\nu$ --- the derivative of the coordinate transformation itself.  These two
transformations are inverses of each other: the contravariant components
$V^\mu$ transform with the Jacobian $\partial x'^\mu / \partial x^\nu$, while
the covariant components $\omega_\mu$ transform with the inverse Jacobian
$\partial x^\nu / \partial x'^\mu$.  This is precisely the transformation law
that defines ``contravariant'' and ``covariant'' in the first place.  The
scalar product $\omega_\mu V^\mu$ is invariant because the two
transformations cancel, which is the coordinate-free content of the summation
convention.
\physnote{The Jacobian matrix $\partial x'^\mu / \partial x^\nu$ is the
derivative (pushforward) of the transition map (\cref{eq:transition-map}).  It
is a linear map between tangent spaces, not the transition map itself.  When
the transition map is a Lorentz transformation, its derivative is the same
Lorentz matrix, and we recover the transformation laws from
\cref{sec:sr-lorentz}.}

As a concrete example, consider polar coordinates $(r, \theta)$ related to
Cartesian $(x, y)$ by $x = r\cos\theta$, $y = r\sin\theta$.  The Jacobian
$\partial x^i / \partial x'^j$ has entries $\partial x / \partial r =
\cos\theta$, $\partial x / \partial \theta = -r\sin\theta$,
$\partial y / \partial r = \sin\theta$,
$\partial y / \partial \theta = r\cos\theta$.  A vector $V$ with Cartesian
components $V^x = 1$, $V^y = 0$ (pointing along the $x$-axis) has polar
components $V^r = \cos\theta$, $V^\theta = -\sin\theta / r$ --- the components
change, but the vector itself is the same geometric arrow.

\paragraph{From components to geometry.}

The tangent-space framework resolves a tension left open from
\cref{ch:special-relativity}.  There, we introduced contravariant and
covariant components as ``upper'' and ``lower'' indices, with transformation
laws stated as rules.  Now we see the geometry behind those rules: upper-index
objects $V^\mu$ are components of tangent vectors, lower-index objects
$\omega_\mu$ are components of one-forms, and the two live in genuinely
different vector spaces.  The metric provides a canonical isomorphism between
them --- the ``translation dictionary'' of \cref{sec:sr-minkowski} --- but
only when the metric is non-degenerate (as every physical spacetime metric is).
Without the metric, the two spaces are distinct.
\implnote{The codebase enforces this distinction at compile time.  The type
\texttt{Vec4 'Contravariant} represents a tangent vector; \texttt{Vec4
'Covariant} represents a one-form.  The phantom \texttt{Variance} parameter
prevents mixing them without an explicit raising or lowering operation.}
\coderef{Math.Tensor.Index}{Variance}

The next section introduces the object that connects these two spaces: the
metric tensor, which generalises the Minkowski metric $\eta_{\mu\nu}$ from
flat spacetime to arbitrary curved manifolds.
